\documentclass[a4paper,10pt]{article}
\usepackage[utf8]{inputenc}
\usepackage[margin=1.5cm, bottom=2cm]{geometry}
\usepackage{fancyhdr}
\usepackage{enumitem}
\usepackage{xcolor}
\usepackage{makeidx}
\usepackage[most]{tcolorbox}
\newtcolorbox{osservazioni}{
  colback=gray!8,    % sfondo grigio chiaro
  colframe=white,      % nessun bordo
  arc=2mm,            % angoli smussati
  boxrule=0pt,        % spessore bordo 0
  fontupper=\footnotesize,
  left=2mm, right=2mm, top=2mm, bottom=2mm,
  before skip=7pt,   % niente spazio prima
  after skip=15pt     % niente spazio dopo
}
\newtcolorbox{osservazioni2}{
  colback=gray!8,    % sfondo grigio chiaro
  colframe=white,      % nessun bordo
  arc=2mm,            % angoli smussati
  boxrule=0pt,        % spessore bordo 0
  fontupper=\small,
  left=5mm, right=5mm, top=3mm, bottom=3mm,
  before skip=0pt,   % niente spazio prima
  after skip=30pt     % niente spazio dopo
}
\pagestyle{fancy}
\usepackage{tikz}
\usetikzlibrary{automata, positioning}
\pagestyle{fancy}

\usepackage{makeidx}

%%%%%%%%%%%%%%%%%%%
\newcommand{\EsercitazioneNum}{07}
\newcommand{\Data}{23/10/2025}
\newcommand{\DocType}{Esercitazione} % o "Eserciziario"
\newcommand{\FullTitle}{\DocType\ \EsercitazioneNum}

\newcommand{\Header}{
\noindent
  \small \textbf{Computabilità, Complessità e Logica - a.a. 2025/26} \hfill \textbf{\FullTitle} \\[0.2cm]
  \scriptsize Matteo Liotta, \texttt{matteo.liotta@studenti.units.it} \hfill \Data \\[0.2cm]
  \noindent\makebox[\linewidth]{\rule{\linewidth}{0.4pt}} \\[0.3cm]}

\newcommand{\NewPageHeader}{
  \vfill
  \noindent\makebox[\linewidth]{\rule{\linewidth}{0.4pt}}
  \newpage
  {\Header}
}
%%%%%%%%%%%%%%%%%%%

% Numero di pagina in basso al centro
\fancyhf{}
\fancyfoot[C] \thepage
\renewcommand{\footrulewidth}{0pt} % niente linea in basso
\renewcommand{\headrulewidth}{0pt} % niente linea in alto

\begin{document}

\footnotesize  % font generale più piccolo
\Header\\[0.7cm]


%%%%%%%%%%%%%%%% FOGLIO 2


%% INDEX
\addcontentsline{toc}{section}{Esercitazione 07: \textit{Macchine di Turing}}

%%
\noindent\Large \textbf{Esercitazione 07}: \textit{Macchine di Turing}\\[0cm]
%\\[0.4cm]i

\begin{osservazioni}
    \textbf{\textit{Nota}}. Sono indicati in \textcolor{blue}{\textit{blu}} gli esercizi che non saranno trattati durante l'esercitazione, bensì lasciati come strumento aggiuntivo di esercitazione individuale. Resta comunque assoluta disponibilità per la loro risoluzione negli incontri successivi, qualora necessario.
\end{osservazioni}

{\normalsize
\begin{itemize}[label=, leftmargin=0pt]

    %% INDEX
    \addcontentsline{toc}{subsection}{Esercizio 47.} 
    %%
    \item \textbf{Esercizio 47 (*)}\\ Sia $L$ linguaggio regolare su $\Sigma=\{a,b\}$, tale per cui 
    {\begin{center}
        $L = \{ w \in \Sigma^* \mid w = a^nb^n,\ \ \  n\geq0 \}$
    \end{center}} 
    Si rappresenti la macchina di Turing tale da riconoscere il linguaggio $L$.\\[0.1cm]

    
    %% INDEX
    \addcontentsline{toc}{subsection}{Esercizio 48.} 
    %%
    \item \textbf{Esercizio 48. (**)}\\ Sia $L$ linguaggio su $\Sigma=\{a,b,c\}$, tale per cui 
    {\begin{center}
        $L = \{ w \in \Sigma^* \mid w = a^{n_1}b^{n_2}c^{n_3}, n_3 = n_1+n_2, \ \ \ n_1,n_2\geq1\}$
    \end{center}} 
    Si rappresenti la macchina di Turing tale da riconoscere il linguaggio $L$.\\[0.1cm]
    

    %% INDEX
    \addcontentsline{toc}{subsection}{Esercizio 49.} 
    %%
    \item \textbf{Esercizio 49. (***)}\\ Si scriva una macchina di Turing $M$ che avendo una stringa $w$ in $\Sigma=\{0, 1\}^*$ termina lasciando sul nastro la stringa $w\# w'$ dove $w'$ è la rappresentazione binaria della lunghezza di $w$.\\[0.1cm]

    
    %% INDEX
    \addcontentsline{toc}{subsection}{Esercizio 50.} 
    %%
    \item \textcolor{blue}{\textbf{Esercizio 50. (*)}}\\ Sia $L$ linguaggio su $\Sigma=\{a,b,c\}$, tale per cui 
    {\begin{center}
        $L = \{ w \in \Sigma^* \mid w = a^{n_1}b^{n_2}c^{n_3}, n_3 = n_1\times n_2, \ \ \ n_1,n_2\geq 0\}$
    \end{center}} 
    Si rappresenti la macchina di Turing tale da riconoscere il linguaggio $L$.\\[0.1cm]

    

    
    
    %% INDEX
    \addcontentsline{toc}{subsection}{Esercizio 51.} 
    %%
    \item \textbf{Esercizio 51. (**)}\\ Si descriva ad alto livello una macchina di Turing $M$ che riceve in ingresso la codifica di altre due macchine di Turing (deterministiche) $M_1$ e $M_2$ e una parola $w$. Tale macchina $M$ deve accettare l'input che riceve sul nastro se la macchina  $M_1$ accetta $w$ in un numero inferiore di passi rispetto a $M_2$.
\end{itemize}
\vfill
\noindent\makebox[\linewidth]{\rule{\linewidth}{0.1pt}}
\end{document}
